\documentclass[letterpaper,10pt]{article}

\usepackage[margin=0.5in]{geometry}
\usepackage[hidelinks]{hyperref}
\usepackage{enumitem}
\usepackage{array}
\usepackage{graphicx}

\pagenumbering{gobble}
\setlength{\parindent}{0pt}
\setlength{\parskip}{0pt}

\setlist[itemize]{leftmargin=*, itemsep=2pt, topsep=2pt, parsep=0pt}

\newcommand{\resumesection}[1]{%
    \vspace{6pt}
    \noindent{\large\bfseries #1}\par
    \vspace{2pt}
    \hrule
    \vspace{4pt}
}

\begin{document}

\begin{center}
    {\LARGE \textbf{Jairaj Sharma}}\\[2.5pt]
    \resizebox{\textwidth}{!}{\mbox{US Citizen \textbar\ Email: \href{mailto:sharma.jairajn@gmail.com}{sharma.jairajn@gmail.com} \textbar\ Phone: 603-921-1481 \textbar\ LinkedIn: \href{https://linkedin.com/in/jairajsharma}{linkedin.com/in/jairajsharma} \textbar\ GitHub: \href{https://github.com/jairaj-S}{github.com/jairaj-S}}}
\end{center}

\resumesection{Summary}
Computer Science graduate with research experience in autonomous systems, machine learning, and defense technology. Skilled at analyzing complex software systems, comparing technical approaches against requirements, and documenting findings clearly. Seeking a Computer Science Patent Examiner position to apply technical analysis and written communication skills to patentability determinations.

\resumesection{Education}
\textbf{Virginia Tech} \hfill Blacksburg, VA\\
Bachelor of Science in Computer Science, GPA: 3.46/4.0 \hfill Aug. 2023 -- May 2026\\
\textit{Relevant Coursework:} Operating Systems, Computer Organization, Programming Languages, Introduction to AI, Data Analytics \& Visualization, Data Structures \& Algorithms

\resumesection{Projects}

\textbf{CanvasInsight} \textbar\ \textit{Flask, React, SQLite, OAuth 2.0, Canvas LTI, LLM API} \hfill Jan. 2026 -- May 2026
\begin{itemize}
    \item Developed and deployed an instructor-facing Canvas LMS tool to deliver automated analysis for student discussion threads, currently utilized in two Summer 2026 Computer Science Ethics courses.
    \item Built automated analysis of large volumes of unstructured text (up to 200 student replies per thread) using the school's LLM capabilities to surface engagement patterns, common misconceptions, and rubric-aligned insights.
    \item Generated organized participation metrics, cross-thread trend summaries, and actionable teaching recommendations directly from structured and unstructured data to significantly reduce manual instructor analysis.
\end{itemize}

\textbf{Patient Data Retrieval Tool @ VTHacks 13 Winner} \textbar\ \textit{JS, Python, Expo, FastAPI} \hfill Sept. 2025
\begin{itemize}
    \item Built an Expo/React Native EMS app for patient identification that captures facial images and surfaces critical medical records, designed to reduce EMT response times.
    \item Designed and prototyped a FastAPI backend supporting facial-recognition matching for dispatcher-to-EMT communication.
\end{itemize}

\resumesection{Experience}

\textbf{Undergraduate Research Assistant} \hfill Sept. 2025 -- June 2026 \textbar\ 4--6 hrs/week\\
\textit{Virginia Tech National Security Institute} \hfill Blacksburg, VA
\begin{itemize}
    \item Engineered 2 programmatic UAS (Unmanned Aerial Systems) scenarios in a Virtual Testing Environment, emulating drone base-incursion events via Python API to support replicable adversarial simulation.
    \item Automated waypoint generation and telemetry capture (position, RF path loss, success metrics) for ground and aerial vehicles, enabling reproducible batch experiments across parameterized configurations.
    \item Designed reproducible experiments using a parameterized generator (PyYAML, SciPy, Latin Hypercube Sampling), ensuring consistent and verifiable results across trials.
    \item Analyzed and documented trial outcomes using pandas and Matplotlib dashboards to support research conclusions.
\end{itemize}

\vspace{4pt}

\textbf{Student Concept Developer, Hacking for Defense Fall 2025 Cohort} \hfill Sept. 2025 -- Dec. 2025 \textbar\ 10 hrs/week\\
\textit{Virginia Tech} \hfill Blacksburg, VA
\begin{itemize}
    \item Selected as 1 of 5 students to lead concept development for single-pilot multi-drone swarm control, translating mission needs into system requirements and risk constraints.
    \item Led 70+ structured interviews with Air Force leadership on base-incursion scenarios and translated findings into written system requirements and evaluation metrics (pilot workload, link latency).
    \item Evaluated and compared drone command-and-control software stacks (Betaflight vs.\ MAVLink/QGroundControl) against FAA and DoD deployability requirements, identifying technical trade-offs and constraints.
\end{itemize}

\vspace{4pt}

\textbf{AI Solutions and Government Bids Intern} \hfill June 2025 -- Sept. 2025 \textbar\ 30 hrs/week\\
\textit{Trinity Bridge Consulting Services} \hfill Nashua, NH
\begin{itemize}
    \item Wrote standardized technical documentation for 9 ML prototypes, describing data collection, model training, and monitoring workflows for government proposal submissions.
    \item Built a National Health and Nutrition Examination Survey 2021-2023 machine learning pipeline to predict dietary supplement ingredient categories from lab and questionnaire data, using Decision Trees, GridSearchCV cross-validation, class balancing, and confusion-matrix analysis.
\end{itemize}

% Keep this section together on one page
\noindent\begin{minipage}{\textwidth}
\resumesection{Technical Skills and Certifications}
\textbf{Certifications:} AWS Certified AI Practitioner (July 2026), AWS Certified Cloud Practitioner (Aug. 2026) \\
\textbf{Languages:} Python, C, Java, JavaScript, SQL, HTML/CSS \\
\textbf{Frameworks/Tools:} AWS Bedrock, AWS Sagemaker, Visual Studio, FastAPI, Flask, React Native, Expo, Docker, \\
Git/GitHub/Gitlab, Linux, REST APIs, YAML, JSON, OpenAI API \\
\textbf{Libraries:} Pandas, NumPy, Matplotlib, SciPy, Scikit-learn \\
\textbf{Drone/UAS Tools:} MAVLink, QGroundControl \\
\textbf{Technical Areas:} Machine learning, autonomous/unmanned systems, software architecture
\end{minipage}

\end{document}
